%
% chapter.tex -- Analytische Funktionen
%
% (c) 2025 Prof Dr Andreas Müller
%
% !TeX spellcheck = de_CH
\chapter{Integration rationaler Funktionen
\label{chapter:ratint}}
\kopflinks{Integration rationaler Funktionen}
Schon früh in der Geschichte der Analysis wurde klar, dass es für
rationale Funktionen immer möglich ist, eine Stammfunktion anzugeben.
Die Methode von Bernoulli, die in
Abschnitt~\ref{buch:ratint:section:bernoulli} dargestellt wird,
findet eine solche Stammfunktion.
Er basiert jedoch auf einer Faktorisierung des Nenners in Linearfaktoren,
deren Existenz der Fundamentalsatz der Algebra zwar garantiert, aber
nicht finden kann.
Die Methode von Bernoulli ist daher im Allgemeinen für ein
Computeralgebrasystem (CAS) nicht durchführbar.
Das Prinzip ist aber verallgemeinerungsfähig, es muss nur eine geeignete
Alternative zur klassischen Partialbruchzerlegung gefunden werden.
Die quadratfreie Faktorisierung einer rationialen Funktion setzt jedoch
eine effiziente Methode voraus, gemeinsame Teiler von Polynomen
zu finden.
Der verallgemeinerte euklidische Algorithmus für Polynome wird in
Abschnitt~\ref{buch:ratint:section:euklid} dargestellt.
Mit der quadratfreien Faktorisierung können somit ein CAS-taugliche
Integrationsalgorithmen konstruiert wird.
Wir stellen die Methode von Hermite in
Abschnitt~\ref{buch:ratint:section:hermite} und die Methode von Horowitz in
Abschnitt~\ref{buch:ratint:section:horowitz} dar.

%
% 1-bernoulli.tex
%
% (c) 2026 Prof Dr Andreas Müller
%
\section{Integration rationaler Funktionen nach Bernoulli
\label{buch:ratint:section:bernoulli}}
\kopfrechts{Integration rationaler Funktionen nach Bernoulli}%
Die Aufgabe, eine Stammfunktion zu finden, gehört zu den zentralen 
Themen jeder Analysis-Grundvorlesung.
Eine besondere Rolle spielt darin der Fall einer rationalen
Funktion als Integrand.
Für diesen Fall gibt es eine vollständige Lösung, mindestens
im Prinzip.
Dieser Abschnitt zeigt diese Lösung von Bernoulli und hebt vor
allem hervor, warum sie für ein Computeralgebrasystem nicht
geeignet ist.

%
% Problemstellung
%
\subsection{Problemstellung}
Die von Bernoulli gelöste Aufgabe ist die Stammfunktion einer
rationalen Funktion zu finden, für die ausserdem angenommen wird,
dass eine Faktorisierung des Nenners in Linearfaktoren bekannt ist.
Der Fundamentalsatz der Algebra garantiert, dass eine solche
Faktorisierung immer exisiert, wenngleich er keine Methode
bereitstellt, die Faktorisierung zu finden.

\begin{aufgabe}
Gegeben ist eine rationale Funktion $f\in \mathbb{C}(z)$ mit
komplexen Koeffizienten.
Finde eine Stammfunktion $F$.
\end{aufgabe}

Aufgrund seiner Annahmen kann Bernoulli annehmen, dass zur
gegebenen rationale Funktion eine Partialbruchzerlegung 
wie in Abschnitt~\ref{buch:analytisch:rational:subsection:partialbruch}
gefunden werden kann.
Der Integrand kann daher als Summe
\begin{align*}
f(z)
=
p(z)
+
\sum_{k=1}^n\sum_{l=1}^{n_k}\frac{A_k^l}{(z-\alpha_k)^l}
\end{align*}
geschrieben werden, wobei $p(z)$ ein Polynom und $\alpha_k$ die Nullstellen
des Nenners von $f(z)$ sind. 
Um eine Stammfunktion zu finden, müssen Stammfunktionen der
einzelnen Summanden gefunden werden.
Der Polynomterm ist dabei das kleinste Problem, da die einzelnen
Summanden $z^n$ die Stammfunktion $z^{n+1}/(n+1)$ haben.
Etwas mehr Arbeit geben die Brüche.

%
% Stammfunktionen für Partialbrüche
%
\subsection{Stammfunktionen für Partialbrüche}
Dank der Partialbruchzerlegung zerfällt in Summanden der Form
\[
\frac{A}{(z-\alpha)^l}
=
A(z-\alpha)^{-l}.
\]
Die Stammfunktion ist
\[
\int \frac{A}{(z-\alpha)^l} \,dz
=
\begin{cases}
\displaystyle
-\frac{A}{(l-1)(z-\alpha)^{l-1}}
&\qquad\text{für $l>1$}
\\[10pt]
A\log(z-\alpha)
&\qquad\text{für $l=1$.}
\end{cases}
\]
Diese Stammfunktionen erlauben eine Stammfunktion für die
Funktion $f(z)$ zu finden.
Sie besteht aus zwei Teilen.
Ein Teil hat die Form einer rationalen Funktion, die aus den
Beiträgen des Polynoms $p(z)$ und den Beiträgen der
Bruchterme $A_k^l/(z-\alpha_k)^l$ mit $l>1$ gelten.
Dazu kommen als zweiter Teil die Funktionen
\[
\sum_{k=1}^n A_k^1\log(z-\alpha_k).
\]
Um eine Stammfunktion eines Integranden im Körper der rationalen Funktionen
zu konstruieren, muss man den Funktionenkörper um Logarithmen der
Linearfaktoren des Nenners erweitern.

%
% Rationale Funktionen mit reellen Koeffizienten
%
\subsection{Rationale Funktionen mit reellen Koeffizienten}
Der Fundamentalsatz der Algebra garantiert eine Faktorisierung 
eines Polynoms mit komplexen Koeffizienten in Linearfaktoren
der Form $z-\alpha$, wobei $\alpha$ eine Nullstelle des Polynoms
ist.
Dies ist jedoch nur möglich, wenn wenn komplexe Nullstellen zugelassen
werden.
Es ist jedoch einfach Polynome anzugeben, die über den reellen Zahlen
nicht in Linearfaktoren zerlegt werden können.
Das Polynom $p(x)=x^2+1$ hat keine reellen Nullstellen.
Über den reellen Zahlen ist $p(x)$ ein irreduzibles Polynom.
Über den komplexen Zahlen ist jedoch $p(x)=(x+i)(x-i)$ eine
Faktorisierung in Linearfaktoren.

%
% Faktorisierung in lineare und quadratische Faktoren
%
\subsubsection{Faktorisierung in lineare und quadratische Faktoren}
Für ein reelles Polynom 
$p(z) = a_0 + a_1z + \dots + a_{n-1}z^{n-1} + a_nz^n$
kann man also eine Faktorisierung in Linearfaktoren
nicht erwarten.
Betrachtet man das Polynom jedoch als ein Polynom mit komplexen
Koeffizienten, dann garantiert der Fundamentalsatz der Algebra eine
Faktorisierung
\[
p(z)
=
a_n
(z-\alpha_1) 
(z-\alpha_2) 
\cdot
\ldots
\cdot
(z-\alpha_n),
\]
wobei die $\alpha_k$ Nullstellen des Polynoms sind.
Es gilt also $p(\alpha_k)=0$ für alle $k$.

Das Polynom $p(x)$ hat reelle Koeffizienten, die $\bar{a}_k=a_k$
für alle $k$ erfüllen.
Daher kann man die Bedinung $p(\alpha_k)=0$ komplex konjugieren
und bekommt
\begin{align*}
0
=
\overline{p(\alpha_k)}
&=
\bar{a}_0
+
\bar{a}_1 \bar{\alpha}_k
+
\bar{a}_2 \bar{\alpha}_k^2
+
\ldots
+
\bar{a}_{n-1} \bar{\alpha}_k^{n-1}
+
\bar{a}_{n} \bar{\alpha}_k^{n}
\\
&=
a_0
+
a_1 \bar{\alpha}_k
+
a_2 \bar{\alpha}_k^2
+
\ldots
+
a_{n-1} \bar{\alpha}_k^{n-1}
+
a_{n} \bar{\alpha}_k^{n}
\\
&=
p(\bar{\alpha}_k).
\end{align*}
Es folgt, dass $\bar{\alpha}_k$ ebenfalls eine Nullstelle ist.
Falls $\alpha_k=\bar{\alpha}_k$ ist, ist $\alpha_k$ eine reelle
Nullstelle.
Falls $\alpha_k\ne \bar{\alpha}_k$ ist, dann ist $\bar{\alpha}_k$
eine der anderen Nullstellen.
Die Nullstellen eines reellen Polynoms sind als entweder reell,
oder wenn sie nicht reell sind, dann treten sie in Paaren von
konjugiert komplexen Nullstellen auf.

Ein paar von konjugiert komplexen Nullstellen führt auf den Faktor
\begin{align}
(z-\alpha)(z-\bar{\alpha})
&=
z^2 -(\alpha+\bar{\alpha}) z + \alpha\bar{\alpha}
\label{buch:ratint:bernoulli:reell:faktoren}
\intertext{Die Summe konjugiert komplexer Zahlen ist
$(a+bi)+(a-bi) = 2a $, also das Doppelte des Realteils.
Daher ist der zu berechnende Faktor}
&=
z^2 -2\operatorname{Re}(\alpha)\, z + |\alpha|^2
\notag
\end{align}
Sowohl $\operatorname{Re}\alpha$ und $|\alpha|^2$ sind reell.
Indem man die Paare konjugiert komplexer Zahlen zu quadratischen
Faktoren wie in \eqref{buch:ratint:bernoulli:reell:faktoren}
zusammenfasst, findet man eine Faktorisierung des Polynoms in
der Form
\begin{equation}
p(z)
=
a_n
(z-\alpha_1)
\cdot\ldots\cdot
(z-\alpha_m)
\cdot
(z^2 + \beta_1 z + \gamma_1)
\cdot \ldots \cdot
(z^2 + \beta_l z + \gamma_l),
\label{buch:ratint:bernoulli:eqn:reellefaktorisierung}
\end{equation}
wobei die $\alpha_1,\dots,\alpha_m$ die reellen Nullstellen sind
und die Faktoren $z^2 +\beta_k z +\gamma_k$ Paare von
konjugiert komplexen, aber nicht reellen Lösungen haben.
Damit die rechte Seite den richtigen Grad hat, muss $n=m+2l$ sein.
\label{buch:ratint:bernoulli:eqn:reellefaktorisierung}

%
% Reelle Partialbruchzerlegung
%
\subsubsection{Reelle Partialbruchzerlegung}
In der Zerlegung
\eqref{buch:ratint:bernoulli:eqn:reellefaktorisierung}
können einzelne Faktoren auch mehrfach vorkommen.
Die Konstruktion der Partialbruchzerlegung produziert
daher für jeden Faktor Partialnenner in Potenzen nicht
grösser als in der Faktorisierung.
Dies ist im folgenden Satz zusammengefasst.

\begin{satz}
\label{buch:ratint:bernoulli:satz:reelle-partialbruchzerlegung}
Ist der Nenner $q(z)$ der rationalen Funktion $f(z)=p(z)/q(z)$ 
ein Polynom mit einer Faktorisierung der Form
\begin{equation}
q(z)
=
a_n
(z-\alpha_1)^{s_1}
\cdots
(z-\alpha_m)^{s_m}
(z^2+\beta_1 z + \gamma_1)^{r_1}
\cdots
(z^2+\beta_l z + \gamma_l)^{r_l}
\end{equation}
dann gibt es eine eindeutig bestimmte Partialbruchzerlegung der Form
\begin{equation}
\renewcommand{\arraycolsep}{2pt}
%\renewcommand{\arraystretch}{1.7}
\begin{array}{rcclcccl}
f(z)
&=&&
\displaystyle
\frac{A_1^1    }{ z-\alpha_1       }
&+& \ldots &+&
\displaystyle
\frac{A_1^{r_1}}{(z-\alpha_1)^{s_1}}
\\
&&&\hspace*{11pt}\vdots&&&&\hspace*{17pt}\vdots
\\[2pt]
&&+&
\displaystyle
\frac{A_m^{1}  }{ z-\alpha_m       }
&+& \ldots &+&
\displaystyle
\frac{A_m^{s_m}}{(z-\alpha_m)^{s_m}}
\\[10pt]
&&+&
\displaystyle
\frac{B_1^1    z + C_1^1    }{ z+\beta_1 z + \gamma_1       }
&+& \ldots &+&
\displaystyle
\frac{B_1^{r_1}z + C_1^{r_1}}{(z+\beta_1 z + \gamma_1)^{r_1}}
\\
&&&\hspace*{24pt}\vdots&&&&\hspace*{29pt}\vdots
\\[2pt]
&&+&
\displaystyle
\frac{B_l^1    z + C_1^l    }{ z+\beta_l z + \gamma_l       }
&+& \ldots &+&
\displaystyle
\frac{B_l^{r_l}z + C_l^{r_1}}{(z+\beta_l z + \gamma_l)^{r_l}}
\end{array}
\label{buch:ratint:bernoulli:eqn:reelle-partialbruchzerlegung}
\end{equation}
Darin sind die Zahlen $A_i^k$, $B_i^k$ und $C_i^k$ reelle Zahlen.
\end{satz}

\begin{proof}
Die Form der Partialbruchzerlegung ist aus der früheren Diskussion,
die zu Satz~\ref{buch:analytisch:ratioinal:satz:partialbruchzerlegung}
geführt hat, klar.
Es muss höchstens noch gezeigt werden, dass diese Zerlegung auch
eindeutig ist.

Offenbar sind $s_i$ Zahlen $A_i^k$ zu bestimmen, insgesamt also
$s_1+\dots+s_m$ Unbekannte für Terme zu den Linearfaktoren.
Zu den quadratischen Faktoren sind je $s_i$ Unbekannte
$B_i^k$ und $C_i^k$ zu bestimmen, also insgesamt
$2r_1+\dots+2r_l$ Unbekannte.
Die Gesamtzahl der Unbekannten ist $s_1+\dots+s_m+2r_1+\dots+2r_l=n$.


Multipliziert man die Gleichung
\eqref{buch:ratint:bernoulli:eqn:reelle-partialbruchzerlegung}
mit dem Nenner $q(z)$, dann bleibt auf der linken Seite das Polynom
$p(z)$ und auf der rechten ein Polynom, dessen Koeffizienten
Linearkombinationen der Unbekannten $A_i^k$, $B_i^k$ und $C_i^k$
sind.
Das Zählerpolynom hat Grad $<n$, es enthält also höchstens $n$
Koeffizienten.
Koeffizientenvergleich liefert dann insgesamt $n$ lineare Gleichungen
für die Unbekannten.
Das entstehende lineare Gleichungssystem hat also $n$ Gleichungen für
$n$ Unbekannte.
\end{proof}

Die reelle Partialbruchzerlegung erlaubt die Berechnung einer
Stammfunktion durch Integration der einzelnen Terme.
Für die linearen Terme ist bereits die Stammfunktion
\[
\int\frac{dz}{z-\alpha}
=
\log (z-\alpha)
\]
bekannt.
Die Stammfunktion für die Terme zu den quadratischen Faktoren
des Nenners wird im nächsten Abschnitt bestimmt.

%
% Stammfunktionen für quadratische Faktoren
%
\subsubsection{Stammfunktionen für quadratische Faktoren}
Wir betrachten jetzt einen quadratischen Nenner der Form
\[
z^2 + \beta z + \gamma
=
(z-\alpha)(z-\bar{\alpha}),
\]
wobei $\beta = 2\operatorname{Re}\alpha$ und $\gamma=|\alpha|^2$ ist.

Als erstes berechnen wir das Integral eines Terms, in dessen Nenner
der quadratische Faktor einmal vorkommt, mit Hilfe der komplexen
Partialbruchzerlegung, die die Form
\begin{align*}
\frac{Bz+C}{z^2+\beta z+\gamma}
&=
\frac{A_1}{z-\alpha}
+
\frac{A_2}{z-\bar{\alpha}}.
\intertext{Durch Multiplizieren mit dem Nenner $(z-\alpha)(z-\bar{\alpha})$
entsteht}
Bz+C
&=
A_1(z-\bar{\alpha})
+
A_2(z-\alpha)
=
(A_1+A_2)z-(A_1\bar{\alpha}+A_2\alpha).
\end{align*}
Durch Koeffizientenvergleich entsteht das lineare Gleichungssystem
\[
\renewcommand{\arraycolsep}{2pt}
\begin{array}{rcrcr}
             A_1 &+&        A_2 &=&  B\\
\bar{\alpha} A_1 &+& \alpha A_2 &=& -C\rlap{,}
\end{array}
\]
Das Gleichungssystem ist am einfachsten mit der Carmerschen Regel
zu lösen:
\begin{align*}
A_1
&=
\frac{
\left|\,\begin{matrix*}[r]
     B       &    1   \\
    -C       & \alpha
\end{matrix*}\,\right|
}{
\left|\,\begin{matrix*}[r]
     1       &    1   \\
\bar{\alpha} & \alpha
\end{matrix*}\,\right|
}
=
\frac{B\alpha+C}{2\operatorname{Im}\alpha}
&
A_2
&=
\frac{
\left|\,\begin{matrix*}[r]
     1       &    B   \\
\bar{\alpha} &   -C  
\end{matrix*}\,\right|
}{
\left|\,\begin{matrix*}[r]
     1       &    1   \\
\bar{\alpha} & \alpha
\end{matrix*}\,\right|
}
=
-
\frac{B\bar{\alpha}+C}{2\operatorname{Im}\alpha}.
\end{align*}
Wir wissen bereits, dass der Realteil von $\alpha$ durch
$\operatorname{Re}\alpha=\frac12 \beta$ ausgedrückt werden kann.
Wegen $\gamma=|\alpha|^2$ und $\beta=2\alpha$ kann man
auch den Imaginärteil von $\alpha$ durch $\beta$ und $\gamma$ ausdrücken,
es ist nämlich
\begin{align*}
(\operatorname{Im}\alpha)^2 &= |\alpha|^2 - (\operatorname{Re}\alpha)^2
\\
\operatorname{Im}\alpha &= \frac12\!\sqrt{4\gamma^2 - \beta^2}.
\end{align*}
Damit kann die Lösung auch
\begin{align*}
A_1
&=
\frac{B\alpha + C}{\!\sqrt{4\gamma^2-\beta^2}}
&&&
A_2
&=
-
\frac{B\bar{\alpha} + C}{\!\sqrt{4\gamma^2-\beta^2}}
\\
&=
\frac{B\operatorname{Re}\alpha + C}{\!\sqrt{4\gamma^2-\beta^2}}
+
\frac{iB\operatorname{Im}\alpha}{2\operatorname{Im}\alpha}
&&&
&=
-
\frac{B\operatorname{Re}\alpha + C}{\!\sqrt{4\gamma^2-\beta^2}}
+
\frac{iB\operatorname{Im}\alpha}{2\operatorname{Im}\alpha}
\\
&=
\frac{B\beta + 2C}{2\!\sqrt{4\gamma^2-\beta^2}}
+
\frac{iB}2
&&&
&=
-
\frac{B\beta + 2C}{2\!\sqrt{4\gamma^2-\beta^2}}
+
\frac{iB}2
\end{align*}

Die Partialbruchzerlegung ist also
\[
\frac{Bz+C}{(z-\alpha)(z-\bar{\alpha})}
=
\frac{1}{2\operatorname{Im}\alpha}
\biggl(
\frac{B\alpha +C}{z-\alpha}
-
\frac{B\bar{\alpha} +C}{z-\bar{\alpha}}
\biggr)
\]


In dieser Form kann die Stammfunktion jetzt mit Logarithmusfunktionen
in der Form
\begin{align*}
\int
\frac{Bz+C}{(z-\alpha)(z-\bar{\alpha})}
\,dz
&=
\frac{1}{2\operatorname{Im}\alpha}
\biggl(
(B\alpha+C)
\log(z-\alpha)
-
(B\bar{\alpha}+C)
\log(z-\bar{\alpha})
\biggr)
\\
&=
\frac{1}{2\operatorname{Im}\alpha}
(B\operatorname{Re}\alpha+C)
\log\frac{z-\alpha}{z-\bar{\alpha}}
\\
&\qquad
+
\frac{1}{2\operatorname{Im}\alpha}
\biggl(
iB\operatorname{Im}\alpha
\log(z-\alpha)
+
iB\operatorname{Im}\alpha
\log(z-\bar{\alpha})
\biggr)
\\
&=
\frac{1}{2\operatorname{Im}\alpha}
(B\operatorname{Re}\alpha+C)
\log\frac{z-\alpha}{z-\bar{\alpha}}
+
\frac{iB}{2}
\log\bigl((z-\alpha)(z-\bar{\alpha})\bigr)
)
\\
&=
\frac{B\operatorname{Re}\alpha+C}{2\operatorname{Im}\alpha}
\log
\frac{z-\alpha}{z-\bar{\alpha}}
+
\frac{iB}2
\log(z^2+\beta z+\gamma)
\end{align*}
dargestellt werden.

%
% Stammfunktion für Potenzen von quadratischen Nennern
%
\subsubsection{Stammfunktion für Potenzen von quadratischen Nennern}
Die Berechnung des vorangegangenen Abschnitts hat gezeigt, wie rationale
Funktionen mit rationalen Nennern integriert werden können.
In der Partialbruchzerlegung treten aber auch noch Summanden mit
Potenzen solcher Nenner auf.
In diesem Abschnitt berechnen wir die Stammfunktionen dieser Terme.
Dabei geht es vor allem um die Vorgehensweise und die dabei 
ausgenützen Eigenschaften der Polynome, die wir später verallgemeinern
werden, wenn allgemeinere Faktorisierungen des Nenners zur Konstruktion
der Partialbruchzerlegung verwendet werden sollen.

In diesem Abschnitt geht es also um die Stammfunktion einer rationalen
Funktion der Form
\begin{equation}
f(z)
=
\frac{Az+B}{(z^2 +\beta z + \gamma)^m}
\qquad\text{mit $m>1$.}
\label{buch:ratint:bernoulli:eqn:quadrpotnenner}
\end{equation}
Wir schreiben auch $q(z) = z^2 +\beta z + \gamma$ und nehmen zudem
an, dass $q(z)$ keine doppelten Nullstellen hat und dass $Az+B$
kein Teiler von $q(z)$ ist.

Zunächst beachten wir, dass der Fall $Az+B=q'(z)$ ein Spezialfall der Form 
\eqref{buch:ratint:bernoulli:eqn:quadrpotnenner}
des Integranden ist.
Für diesen speziellen Zähler kann das Integral sofort ermittelt
werden.

\begin{satz}
\label{buch:ratint:bernoulli:satz:einfacheintegrale}
Unter den formulierten allgemeinen Voraussetzungen an $q(z)$ ist
\begin{equation}
\int
\frac{q'(z)}{q(z)^m}\,dz
=
\begin{cases}
\log q(z)                  &\qquad\text{für $m=1$}\\[5pt]
\displaystyle
\frac{1}{(1-m)q(z)^{m-1}}  &\qquad\text{für $m>1$.}
\end{cases}
\label{buch:ratint:bernoulli:eqn:einfacheintegrale}
\end{equation}
\end{satz}

\begin{proof}
Der Beweis erfolgt durch Nachrechnen der Ableitung
auf der rechten Seite von \eqref{buch:ratint:bernoulli:eqn:einfacheintegrale}.
Für den Fall $m=1$ rechnet man
\begin{align*}
\frac{d}{dz}\log q(z)
&=
\frac{1}{q(z)}\cdot q'(z)
\intertext{nach der Kettenregel.
Für $m>1$ gilt}
\frac{d}{dz}\frac{1}{(1-m)q(z)^{m-1}}
&=
\frac{1}{1-m}
\frac{d}{dz}
q(z)^{1-m}
=
\frac{1}{1-m} (1-m)q(z)^{-m}\cdot q'(z)
=
\frac{q'(z)}{q(z)^m}.
\end{align*}
Damit ist der Satz bewiesen.
\end{proof}

Der Zähler $Az+B$ ist im Allgemeinen von $q'(z)$ verschieden,
trotzdem gibt es eine Verbindung, die später ermöglichen wird,
das Integral zu vereinfachen.

\begin{satz}
\label{buch:ratint:bernoulli:satz:euklidspezial}
Unter den obigen allgemeinen Voraussetzungen gibt es eine Zahl
$s\in\mathbb{R}$ und ein lineares Polynom $t(z)=uz+v\in\mathbb{R}[z]$
derart, dass
\begin{equation}
s\cdot q(z) + t(z)\cdot q'(z) = Az+B.
\label{buch:ratint:bernoulli:eqn:euklidspezial}
\end{equation}
\end{satz}

\begin{proof}
Man setzt die Polynome in die Gleichung
\eqref{buch:ratint:bernoulli:eqn:euklidspezial}
ein und erhält
\begin{align*}
s\cdot(z^2+\beta z+\gamma) + (uz+v)\cdot (2z+\beta) &=  Az+B
\intertext{oder nach Ausmultiplizieren}
sz^2 + \beta s z + \gamma s + 2uz^2 +(\beta u+2v)z +\beta v &= Az + B
\intertext{oder zusammengefasst}
(s+2u) z^2 + (\beta s + \beta u+2v)z +(\beta s + \beta v) &= Az + B.
\end{align*}
Durch Koeffizientenvergleich erhält man das lineare Gleichungssystem
\[
\renewcommand{\arraycolsep}{2pt}
\begin{array}{rcrcrcr}
       s&+&     2 u& &        &=& 0 \\
 \beta s&+& \beta u&+&     2 v&=& A \\
\gamma s& &        &+& \beta v&=& B
\end{array}
\]
mit der Determinanten
\[
\left|\,
\begin{matrix*}[r]
     1 &     2 &     0 \\
 \beta & \beta &     2 \\
\gamma &     0 & \beta
\end{matrix*}
\,
\right|
=
\beta^2 + 4\gamma - 2\beta^2
=
-(\beta^2-4\gamma).
\]
Der Klammerausdruck ist die Diskriminante des quadratischen
Polynoms $q(z)$.
Sie ist von $0$ verschieden, da angenommen wurde, dass $q(z)$
keine doppelten Nullstellen hat.
Die Koeffizientenmatrix ist somit regulär und das Gleichungssystem
hat eine eindeutig bestimmte Lösung, was den Satz beweist.
\end{proof}


\begin{satz}
\label{buch:ratint:bernoulli:satz:reduktion}
Unter den gegebenen allgemeinen Voraussetzungen lässt sich mit den
Faktoren $s$ und $t(z)$ von
Satz~\eqref{buch:ratint:bernoulli:satz:euklidspezial}
das Integral durch ein Integral mit niedrigerem Exponenten $m$ ausdrücken,
nämlich
\begin{equation}
\int
\frac{Az+B}{q(z)^m}
\,dz
=
\int
\frac{s+t'(z)/(m-1)}{q(z)^{m-1}}
\,dz
+
\frac{t(z)}{(1-m)q(z)^{m-1}}
\label{buch:ratint:bernoulli:eqn:reduktion}
\end{equation}
\end{satz}

\begin{proof}
Setzt man die
\eqref{buch:ratint:bernoulli:eqn:euklidspezial}
in den Integranden ein, erhält man
\begin{align*}
\int
\frac{Az+B}{q(z)^m}
\,dz
&=
\int
\frac{s\cdot q(z)}{q(z)^m}
\,dz
+
\int
\frac{t(z) \cdot q'(z)}{q(z)^m}
\,dz.
\intertext{Im ersten Integral kann ein Faktor $q(z)$ weggekürzt werden.
Im zweiten Integral kann partielle Integration angewendet werden,
wobei $q'(z)/q(z)^m$ der Faktor ist, der integriert wird und $t(z)$
der Faktor, der abgeleitet wird.
Für die Stammfunktion von $q'(z)/q(z)^m$ wird
Satz~\eqref{buch:ratint:bernoulli:satz:einfacheintegrale}
verwendet.
So erhält man}
&=
\int
\frac{s}{q(z)^{m-1}}
\,dz
+
t(z)\cdot\frac{1}{(1-m)q(z)^{m-1}}
-
\int
\frac{t'(z)}{(1-m)q(z)^{m-1}}
\,dz.
\\
&=
\int
\frac{s+t'(z)/(m-1)}{q(z)^{m-1}}
\,dz
+
\frac{t(z)}{(1-m)q(z)^{m-1}}.
\end{align*}
Dies ist die behauptete Form des Integrals.
\end{proof}

Der letzte Satz zeigt, dass sich alle Integrale mit Nenner
$q(z)^m$ mit $m>1$ immer auflösen lassen in eine rationale
Funktion und ein Integral mit Nenner $q(z)^{m-1}$.
Wiederholte Anwendung dieser Eigenschaft ermöglicht,
das Integral soweit zu reduzieren, dass am Ende in Integral
mit dem Nenner $q(z)$ übrig bleibt, wie es im vorangegangenen
Abschnitt berechnet wurde.

%
% Die allgemeine Form der Stammfunktion
%
\subsubsection{Die allgemeine Form der Stammfunktion}
Die vorangegangenen Berechnungen zeigen, dass das Integral einer
rationalen Funktion $f(z)$ sich immer als eine Summe von Termen
der folgenden Form darstellen lassen.
\begin{enumerate}
\item
Für jeden Linearfaktor $z-\alpha$ sind rationale Terme der Form
\[
\frac{A}{(z-\alpha)^m}
\]
möglich, wobei $m$ nicht grösser sein kann als die Potenz, mit der
der Faktor im Nenner von $f(z)$ vorkommt.

Ausserdem ist ein Term der Form
\[
A \log (z-\alpha)
\]
möglich.
\item
Für jeden quadratischen Faktor der Form $q(z)=z^2+\beta z+\gamma$
sind rationale Terme der Form $t(z) / q(z)^m$ möglich, wobei $t(z)$
ein lineares Polynom ist und $m$ nicht grösser sein kann als die
Potenz, in der $q(z)$ im Nnner von $f(z)$ vorkommt.

Ausserdem ist ein Term der Form
\[
\log q(z)
=
\log (z^2 +\beta z + \gamma)
\]
möglich.
\end{enumerate}

%
% Die der Methode zugrunde liegenden Prinzipien
%
\subsubsection{Die der Methode zugrunde liegenden Prinzipien}
Es wurde bereits früher festgehalten, dass die dargestellte Methode
für ein Computeralgebrasystem nicht durchführbar ist, weil eine
Methode zur Bestimmung der Faktorisierung des Nenners in Linearfaktoren
oder quadratische reelle Faktoren fehlt.
Es muss daher eine alternative Vorgehensweise gefunden werden,
die mit einer leichter zu findenden Faktorisierung des Nenners
arbeiten kann.
Die Faktorisierung muss immer noch einige Eigenschaften erfüllen,
damit die einzelnen Schritte verallgemeinert werden können.

\begin{enumerate}
\item
Die Partialbruchzerlegung wurde dazu verwendet, die 
rationale Funktion in Term mit speziell einfachen Nennern
aufzuspalten.
In jedem Term kamen höchstens Potenzen eines einzigen Faktors
des Nenners vor.
Wie in der Darstellung der Partialbruchzerlegung angedeutet
wurde, müssen die Faktoren teilerfremd sein.
\item
Für den Reduktionsschritt von
Satz~\ref{buch:ratint:bernoulli:satz:reduktion}
wurde benötigt, dass sich jeder denkbare Zähler in der Form
\[
s(z) \cdot q(z) + t(z)\cdot q'(z)
\]
ausdrücken lässt, wobei $q(z)$ ein einzelner Faktor in der
Nennerfaktorisierung ist.
Dies ist genau dann möglich, wenn $q(z)$ und $q'(z)$ teilerfremd
sind, also keine gemeinsame Nullstelle haben.
Dies wieder ist gleichbedeutend damit, dass $q(z)$ nur einfache
Nullstellen hat.
\item
Es wird ein Algorithmus benötigt, der die Polynom $s(z)$ und
$t(z)$ auf effiziente Art und Weise finden kann.
\item
Es müssen Integrale der Form 
\[
\int \frac{p(z)}{q(z)}\,dz
\]
gefunden werden, wobei $p(z)$ ein beliebiges Polynom ist.
\end{enumerate}

Die Aufgabe von Punkt~3, das Finden einer Darstellung eines
beliebigen Zahlerpolynoms Zwecks Reduktion des Exponenten im
Nenner wird durch den euklidischen Algorithmus ermöglicht,
der im Abschnitt~\ref{buch:ratint:section:euklid} dargestellt
wird.

Eine geeignete Faktorisierung, die immer gefunden werden kann,
ist die sogenannte quadratfreie Faktorisierung, die in Abschnitt 
\ref{buch:ratint:section:hermite}
zusammen mit der Reduktion der Integrale von rationalen
Funktionen mit Potenzen des Nennerfaktors vorgestellt wird.



%
% 2-euklid.tex
%
% (c) 2026 Prof Dr Andreas Müller
%
\section{Der euklidische Algorithmus
\label{buch:ratint:section:euklid}}
\kopfrechts{Der euklidische Algorithmus}%
Die Konstruktion der Partialbruchzerlegung hat erfordert, dass gemeinsame
Faktoren von Polynomen gefunden werden mussten.
Es ist natürlich immer denkbar, eine Faktorisierung der beteiligten
Polynome zu suchen, und die gemeinsamen Faktoren zu identifizieren.
Der Fundamentalsatz der Algebra garantiert zwar, dass solche
Faktoren über $\mathbb{C}$ existieren, für die Anwendung in einem
Computeralgebrasystem brauchen wir aber eine Faktorisierung über
$\mathbb{Q}$ oder einer Körpererweiterung von $\mathbb{Q}$, die in nur
endlich vielen Schritten erfolgt ist.

Das Finden der Faktoren ist eine weitere Schwierigkeit.
Selbst über $\mathbb{C}$ können die vom Fundamentalsatz garantierten
Linearfaktoren für Polynome vom Grad $\ge 5$ algorithmisch im
Allgemeinen nicht gefunden werden.
Es ist daher nicht angebracht, auf dem Umweg über die Faktorisierung
nach gemeinsame Faktoren zu suchen.

Es gibt aber eine direkte Möglichkeit, gemeinsame Faktoren zu finden.
In seiner ursprünglichen Version, die in
Abschnitt~\ref{buch:ratint:euklid:subsection:ganzezahlen}
beschrieben wird, findet der euklidische Algorithmus den grössten
gemeinsamen Teiler $r$ zweier ganzer Zahlen $a$ und $b$.
Die erweiterte Form ist zudem in der Lage, den grössten gemeinsamen
Teiler als $r=sa+tb$ zu schreiben mit ganzen Zahlen $s,t\in\mathbb{Z}$.
Diese Version kann dazu verwendet werden, die multiplikative Inverse 
in einem endlichen Körper zu finden.

Der euklidische Algorithmus ist nicht auf ganze Zahlen beschränkt.
Er lässt sich in jedem sogenannten euklidischen Ring durchführen.
Dazu gehören auch die Polynomringe.
Der grösste gemeinsame Teiler von zwei Polynomen ist ein gemeinsamer
Faktor, der gemeinsame Nullstellen enthält.
Eine mehrfache Nullstelle eines Polynoms $p(X)$ führt zu einem gemeinsamen
Teiler von $p'(X)$ und $p(X)$, auch solche mehrfache Faktoren
können also mit dem euklidischen Algorithmus gefunden werden.

%
% Der euklidische Algorithmus für ganz Zahlen
%
\subsection{Der euklidische Algorithmus für ganze Zahlen
\label{buch:ratint:euklid:subsection:ganzezahlen}}
In diesem Abschnitt wird der euklidische Algorithmus für ganze
Zahlen entwickelt.
Diese Form des Algorithmus ist seit dem Altertum bekannt und wird
oft bereits in der Mittelschule als Beispiel ein sowohl historisch
wie auch mathematisch bedeutsames Beispiel eines Algorithmus
unterrichtet.
Im Abschnitt~\ref{buch:ratint:euklid:subsection:polynome} wird dann
der Algorithmus auf Polynome verallgemeinert.

%
% Der Grundalgorithmus
%
\subsubsection{Der Grundalgorithmus}
Seien $a>b$ natürliche Zahlen und $q$ und $r$ so gewählte
natürliche Zahlen, dass $a=bq+r$ mit $r<b$.
Der grösste gemeinsame Teiler $g=\operatorname{ggT}(a,b)$ teilt dann
sowohl $a$ als auch $b$, er muss daher auch $r$ teilen.
Wiederholt man die Rechnung mit $b$ anstelle von $a$ und $r$ anstelle
von $b$ folgt, dass $g$ erneut den Rest teilt.
Da der Rest immer kleiner wird, gibt es einen Moment, bei dem
er erstmals verschwindet.
Der Divisor $b$ ist dann der grösste gemeinsame Teiler.

\begin{beispiel}
\label{buch:ratint:euklid:bsp:euklidzahlen}
Für die Zahlen $a=36$ und $b=28$ ergibt sich die Rechnung
\begin{center}
\begin{tabular}{|>{$}r<{$}>{$}r<{$}|>{$}r<{$}>{$}r<{$}|}
\hline
 a&  b& q& r\\
\hline
36& 28& 1& 8\\
28&  8& 3& 4\\
 8&  4& 2& 0\\
\hline
\end{tabular},

\end{center}
aus der man $\operatorname{ggT}(36,28)=4$ ablesen kann.
\end{beispiel}

%
% Der erweiterte Algorithmus
%
\subsubsection{Der erweiterte Algorithmus}
Wir analysieren euklidischen Algorithmus nochmals etwas im Detail.
Dazu führen wir die folgende Notation ein.
Wir beginnen mit $a_0=a$ und $b_0=b$.
Wir konstruieren dann $q_0$ und $r_0$ mit $a_0=q_0b_0+r_0$ und 
setzen anschliessend $a_1 = b_0$ und $b_1 = r_0$.
Dieser Prozess wird für $k>0$ gemäss
\begin{align*}
a_k &= b_{k-1} &
b_k &= r_{k-1} &
&\text{und}&
a_k &= q_k b_k + r_k 
&&\Rightarrow&
r_k &= a_k - q_kb_k
\end{align*}
wiederholt, bis $r_{n+1}=0$ ist.
Der grösste gemeinsame Teiler ist dann $r_n$.
Rechnet man ausgehend von $r_n$ rückwärts, findet man
\begin{align*}
r_n &= a_n - q_n b_n
\\
    &= b_{n-1} - q_n r_{n-1}
\\
    &= b_{n-1} - q_n ( a_{n-1} - q_{n-1} b_{n-1} )
     = q_n a_{n-1} + (1 + q_n q_{n-1}) b_{n-1}
\\
    &\phantom{1}\vdots
\intertext{Durch wiederholte Anwendung der Definitionsformeln
für $a_k$ und $b_k$ finden wir schliesslich ganze Zahlen $s$ und $t$
mit der Eigenschaft}
    &= s a_0 + t b_0.
\end{align*}

\begin{beispiel}
Für das Beispiel~\ref{buch:ratint:euklid:bsp:euklidzahlen} folgt
aus der zweiten Zeile der Tabelle
\[
\operatorname{ggT}(a,b)
=
4
= 
28 - 3\cdot 8
\]
Da 8 der Rest der vorangegangenen Zeile ist, kann man ihn durch
$36-1\cdot 28$ ersetzen und erhält
\[
4
=
28 - 3 \cdot (36-28)
=
4\cdot 28 - 3\cdot 36.
\]
Damit sind die Zahlen $s=-3$ und $t=4$ gefunden.
\end{beispiel}

Der erweiterte euklidische Algorithmus kann also nicht nur den 
grössten gemeinsamen Teiler von $a$ und $b$ finden, sondern zusätzlich
Zahlen $s$ und $t$ derart, dass
\[
sa + tb
=
\operatorname{ggT}(a,b)
\]
gilt.
Die Berechnung nach Abschluss ist jedoch eher unzweckmässig.
Gewünscht ist eine Erweiterung des Grundalgorithmus, die die Zahlen
$s$ und $t$ gleich mitberechnet.

\begin{satz}
Gegeben sind die natürlichen Zahlen $a$ und $b$.
Der folgende Algorithmus berechnet den grössten gemeinsamen Teiler $g$
und Zahlen $s$ und $t$ mit $g=sa+tb$.
\begin{enumerate}
\item
Setze $a_0=a$, $b_0=b$, $k=0$ und $s_{-2}=1$, $t_{-2}=0$, $s_{-1}=0$
und $t_{-1}=1$
\item
Finde $q_k$ und $q_k$ mit $a_k= q_kb_k+r_k$.
\item
Berechne
$s_k=s_{k-2}-q_ks_{k-1}$ 
und
$t_k=t_{k-2}-q_kt_{k-1}$.
\end{enumerate}
Für alle $k$ gilt
$r_k = s_ka+t_kb$
Für den grössten Index $n$, für den $r_n\ne 0$ ist, folgt
\[
\operatorname{ggT}(a,b)
=
s_na+t_nb
\qquad\text{und}\qquad
0
=
s_{n+1}a+t_{n+1}b.
\]
\end{satz}

\begin{proof}
Wir schreiben die Division als Matrizenprodukt:
\[
\begin{pmatrix}
b_k\\
r_k
\end{pmatrix}
=
\underbrace{
\begin{pmatrix}
 0 & 1   \\
 1 &-q_k
\end{pmatrix}
}_{\displaystyle = Q(q_k)}
\begin{pmatrix}
a_k\\
b_k
\end{pmatrix}
\]
Wiederholte Anwendung ergibt
\[
\begin{pmatrix}
r_n\\
0
\end{pmatrix}
=
Q(q_{n+1})
Q(q_n)
\cdots
Q(q_1)
Q(q_0)
\begin{pmatrix}
a\\ b
\end{pmatrix}.
\]
Das Produkt der Matrizen enthält in der ersten Zeile die gesuchten
Zahlen $s_n$ und $t_n$ und in der zweiten Zeile die Zahlen
$s_{n+1}$ und $t_{n+1}$.

Da das Matrizenprodukt assoziativ ist, kann das Produkt von rechts
her aufgelöst werden.
Dazu definieren wir die Matrix
\[
S_k
= 
\begin{pmatrix}
s_{k-1} & t_{k-1} \\
s_{k}   & t_{k} 
\end{pmatrix}
\]
mit der initialen Matrix
\[
S_{-1}
=
\begin{pmatrix}
s_{-2} & t_{-2} \\
s_{-1} & t_{-1} 
\end{pmatrix}
=
\begin{pmatrix}
1 & 0\\
0 & 1
\end{pmatrix}.
\]
Die Multiplikation mit $Q(q_k)$ mit $S_{k-1}$ erzeugt $S_k$, es gilt
daher die Rekursionsformel
\begin{align*}
S_k
&= 
\begin{pmatrix}
s_{k-1} & t_{k-1} \\
s_{k}   & t_{k} 
\end{pmatrix}
=
Q(q_k)
S_{k-1}
=
Q(q_k)
\begin{pmatrix}
s_{k-2} & t_{k-2} \\
s_{k-1} & t_{k-1} 
\end{pmatrix}
\\
&=
\begin{pmatrix}
0&1\\
1&-q_k
\end{pmatrix}
\begin{pmatrix}
s_{k-2} & t_{k-2} \\
s_{k-1} & t_{k-1} 
\end{pmatrix}
\\
&=
\begin{pmatrix*}[r]
s_{k-1}            & t_{k-1}           \\
s_{k-2}-q_ks_{k-1} & t_{k-2}-q_kt_{k-1} 
\end{pmatrix*}.
\end{align*}
In der zweiten Zeile stehen die Rekursionsformeln für die
Zahlen $s_k$ und $t_k$.
\end{proof}

\begin{beispiel}
\label{buch:ratint:euklid:bsp:erweitert}
Wir führen den erweiterten euklidischen Algorithmus für $a=43$ und $b=47$
durch:
\begin{center}
\begin{tabular}{|>{$}r<{$}|>{$}r<{$}>{$}r<{$}|>{$}r<{$}>{$}r<{$}|>{$}r<{$}>{$}r<{$}|}
\hline
  k &  a_k &  b_k &  q_k &  r_k &  s_k &  t_k \\
\hline
 -2 &      &      &      &      &    1 &    0 \\
 -1 &      &      &      &      &    0 &    1 \\
  0 &   43 &   47 &    0 &   43 &    1 &    0 \\
  1 &   47 &   43 &    1 &    4 &   -1 &    1 \\
  2 &   43 &    4 &   10 &    3 &   11 &  -10 \\
  3 &    4 &    3 &    1 &    1 &  -12 &   11 \\
  4 &    3 &    1 &    3 &    0 &   47 &  -43 \\
\hline
\end{tabular}
\end{center}
In der letzten Zeile wird der Rest $0$, der grösste gemeinsame Teiler ist
also $\operatorname{ggT}(a,b) = r_3=1$.
In der Zeile $k=3$ kann man die beiden Zahlen $s=s_3=-12$ und $t=t_3=11$
ablesen, und tatsächlich ist
\[
-12\cdot 43 + 11\cdot 47 =  1
\qquad\text{und}\qquad
47\cdot 43 - 43\cdot 47 = 0.
\qedhere
\]
\end{beispiel}

Das Beispiel illustriert auch, dass die im Grundalgorithmus formulierte
Bedingung, dass $a>b$ sein muss, unwichtig ist.
Wenn nämlich $a<b$ ist, dann ist der erste Quotient $q_0=0$ und der erste
Rest $r_0=a$.
Daraus ergibt sich dann $a_1=b$ und $b_1=a$, d.~h.~die beiden Zahlen $a$
und $b$ haben die Plätze getauscht.
Auch die Zahlen $s_0=1$ und $t_0=0$ tauschen die Plätze der Anfangsbedingungen
für die Folgen $s_k$ und $t_k$.

%
% Anwendung des erweiterten Algorithmus
%
\subsubsection{Anwendung des erweiterten Algorithmus}
Der erweiterte euklidische Algorithmus kann dazu verwendet werden,
die multiplikative inverse in einem endlichen Körper zu finden.
Sei $p$ eine Primzahl und $\mathbb{F}_p$ die Menge der Reste modulo $p$.
Die multiplikative Inverse eines Elementes $a\in\mathbb{F}_p$ kann
jetzt wie folgt gefunden werden.
Da $p$ ein Primzahl ist, sind $a$ und $p$ teilerfremd.
Der erweiterte euklidische Algorithmus liefert zwei Zahlen $s$ und $t$
mit
\[
sa + tp = \operatorname{ggT}(a,p) = 1.
\]
Modulo $p$ wird diese Gleichung zu
\[
sa \equiv 1 \mod p.
\]
Der Rest von $s$ modulo $p$ ist daher die multiplikative Inverse $a$.

\begin{beispiel}
Für $p=47$ und $a=43$ ergibt der in
Beispiel~\ref{buch:ratint:euklid:bsp:erweitert}
durchgeführte, erweiterte euklidische Algorithmus

\[
-12 \cdot 43
+
11\cdot 47
= 1.
\]
Die multiplikative Inverse von $43$ im Körper $\mathbb{F}_{47}$ ist
daher $-12\equiv 35\mod 47$.
Tatsächlich ist 
\[
35\cdot 43 = 1505 = 32\cdot 47 + 1,
\]
wie behauptet.
\end{beispiel}


%
% Der euklidische Algorithmus für Polynome
%
\subsection{Der euklidische Algorithmus für Polynome
\label{buch:ratint:euklid:subsection:polynome}}
Der euklidische Algorithmus funktioniert auch für Polynome.
Dies ermöglicht, rationale Funktionen effizient zu kürzen.
Später werden wir den erweiterten euklidischen Algorithmus
auch verwenden, um die quadratfreie Faktorisierung und die 
zugehörige Partialbruchzerlegung einer rationalen Funktion
zu finden und um Integrale von rationalen Funktionen auf 
einfachere Form zu reduzieren.

%
% Polynomdivision
%
\subsubsection{Polynomdivision}
Auch für Polynome gibt es eine Division mit Rest.
Zu zwei Polynomen $a(x)$ und $b(x)$ gibt es Polynome $q(x)$ und $r(x)$,
für die
\[
a(x) = b(x)q(x)+r(x)
\]
gilt.
Ausserdem muss der Grad des Restpolynoms $r(x)$ kleiner sein als
der Grad von $b(x)$: $\deg r(x) < \deg b()$.

\begin{beispiel}
\label{buch:ratint:euklid:bsp:polynomdivision}
Man berechen Quotient und Rest der Polynome
\[
a(x)
=
x^4+x^3-4x^2+x-6
\qquad\text{und}\qquad
b(x)
=
x^2+5x+6.
\]
\[
\renewcommand{\arraycolsep}{1.2pt}
\begin{array}{rcrcrcrcrcrcrcrcrcrcr}
x^4&+&  x^3&-& 5x^2&+&  x&-& 6&:& x^2&+& 5x&+& 6 &=& x^2 &-& 4x&+& 9\\
x^4&+& 5x^3&+& 6x^2& &   & &  & &    & &   & &   & &     & &   & &  \\
\cline{1-5}
   &-& 4x^3&-&11x^2&+&  x& &  & &    & &   & &   & &     & &   & &  \raisebox{4pt}{\mathstrut}\\
   &-& 4x^3&-&20x^2&-&24x& &  & &    & &   & &   & &     & &   & &  \\
\cline{2-7}
   & &     & & 9x^2&+&25x&-& 6& &    & &   & &   & &     & &   & &  \raisebox{4pt}{\mathstrut}\\
   & &     & & 9x^2&+&45x&+&54& &    & &   & &   & &     & &   & &  \\
\cline{5-9}
   & &     & &     &-&20x&-&60& &    & &   & &   & &     & &   & &  \raisebox{3pt}{\mathstrut}\\
\cline{6-9}
\end{array}
\]
Somit ist der Quotient $q(x)=x^2-4x+9$ und der Rest $r(x)=-20x-60$.
\end{beispiel}

Die Division von Polynomen ist die Basis des Algorithmus zur Bestimmung
des grössten gemeinsamen Teilers.
Es ist aber zu beachten, dass die Division einen Rest liefert,
der nicht unbedingt ein normiertes Polynom ist.
Es kann daher nötig sein, den Rest zu normieren.
Im Beispiel ist der normierte Rest $x+3$.

%
% Gemeinsame Teiler für Polynome
%
\subsubsection{Gemeinsame Teiler für Polynome}
Der euklidische Algorithmus für den grössten gemeinsamen Teiler
funktioniert unverändert auch für Polynome.
Der Polynomgrad des Restes nimmt dabei ständig ab, bis der Rest
0 wird.

\begin{beispiel}
Man finde den grössten gemeinsamen Teiler der Polynome
\[
a(x)
=
x^4+x^3-4x^2+x-6
\qquad\text{und}\qquad
b(x)
=
x^2+5x+6.
\]
Dazu müssen die Quotienten und Reste berechnet werden, der
erste wurde bereits im
Beispiel~\ref{buch:ratint:euklid:bsp:polynomdivision}
berechnet:
\begin{center}
\begin{tabular}{|>{$}r<{$}|>{$}l<{$}>{$}l<{$}|>{$}l<{$}>{$}l<{$}|}
\hline
 k & a_k              & b_k      & q_k                & r_k
\raisebox{2.5pt}{\mathstrut}\\[1pt]
\hline
 0 & x^4+x^3-5x^2+x-6 & x^2+5x+6 & x^2-4x+9           & -20x-60
\raisebox{3pt}{\mathstrut}\\
 1 & x^2+5x+6         & -20x-60  & -\frac{1}{20}(x+2) &       0
\\[2pt]
\hline
\end{tabular}
\end{center}
In der zweiten Zeile lässt sich ablesen, dass der Rest $=0$ ist,
daher ist der grösste gemeinsame Teiler $-20x-60$.
Dieses Polynom ist nicht normiert, da die Koeffizienten rational sind,
kann man durch $-20$ teilen und erhält das normierte Polynom $x+3$.
\end{beispiel}

Das Beispiel zeigt auch, dass der gefundene grösste gemeinsame Teiler
nicht ein normiertes Polynom ist.
Da wir nur an Polynomen mit rationalen Koeffizienten interessiert
sind, kann immer durch den Leitkoeffizient geteilt werden.
Arbeitet man dagegen nur mit ganzzahligen Koeffizienten, wird im 
Beispiel bereits der Schritt $k=2$ nicht mehr durchführbar.

%
% Der erweiterte euklidische Algorithmus für Polynome
%
\subsubsection{Der erweiterte euklidische Algorithmus für Polynome}
Auch der erweiterte euklidische Algorithmus lässt sich übertragen.

\begin{beispiel}
Für die Polynome $a(x) = x^3+2x^2-x-2$ und $b(x)=x^2-x-2$
finde man Polynome $s(x)$ und $t(x)$ mit
$\operatorname{ggT}(a,b)=s(x)a(x)+t(x)b(x)$.
\begin{center}
\renewcommand{\arraycolsep}{1.2pt}
\begin{tabular}{|>{$}r<{$}|>{$}l<{$}>{$}l<{$}|>{$}l<{$}>{$}l<{$}|>{$}l<{$}>{$}l<{$}|}
\hline
  k & a_k(x)        & b_k(x)  & q_k(x)       & r_k(x) & s_k(x) & t_k(x) \raisebox{2.5pt}{\mathstrut}\\[1pt]
\hline
 -1 &               &         &              &        &             1 &                0 \\
 -2 &               &         &              &        &             0 &                1 \\
  0 & x^3+2x^2-2x-2 & x^2-x-2 & x+3          & 4x + 4 &             1 &             -x-3 \\
  1 & x^2-x-2       & 4x+4    & \frac14(x-2) & 0      & -\frac14(x-2) & \frac14(x^2+x-2) \\[2pt]
\hline
\end{tabular}
\end{center}
Tatsächlich können wir nachrechnen, dass
\begin{align*}
s_0(x)
\cdot
a(x)
+
t_0(x)
\cdot
b(x)
&=
1
\cdot
(x^3+2x^2-x-2)
-
(x+3)
\cdot
(x^2-x-2)
\\
&=
(x^3+2x^2-x-2)
-
(x^3+2x^2-5x-6)
\\
&=
4x+4
\\
s_1(x)
\cdot
a(x)
+
t_1(x)
\cdot
b(x)
&=
-
\frac14(x-2)
\cdot
(x^3+2x^2-x-2)
+
\frac14(x^2+x-2)
\cdot
(x^2-x-2)
\\
&=
-
\frac14(x^4-5x^2+4)
+
\frac14(x^4-5x^2+4)
\\
&=0.
\end{align*}
Verlangt man ein normiertes Polynom für den grössten gemeinsamen Teiler,
kann man die Koeffizienten in $\mathbb{Q}$ durch $4$ dividieren und 
erhält $\operatorname{ggT}(a,b)=x+1$.
Die Faktoren $s(x)$ und $t(x)$ müssen ebenfalls durch 4 geteilt werden
und es gilt
\[
\frac14
\cdot
(x^3+2x^2-x-2)
-\frac14(x+3)
(x^2-x-2)
=
x+1.
\]

Den grössten gemeinsamen Teiler kann man auch mithilfe der Faktorisierung
\begin{align*}
a(x)
&=
(x-1){\color{darkred}(x+1)}(x+2)
\\
b(x)
&=
(x-2){\color{darkred}(x+1)}
\end{align*}
der Polynome finden.
Der einzige gemeinsame Faktor ist $\operatorname{ggT}(a,b)={\color{darkred}x+1}$.
Aus der Faktorisierung sind jedoch die Faktoren $s(x)$ unt $t(x)$
nicht direkt ablesbar.
\end{beispiel}

%
% Euklidische Ringe
%
\subsubsection{Euklidische Ringe}
Der euklidische Algorithmus benötigt erst mal, dass es keine
pathologischen Faktoren gibt, die ein Produkt auslöschen können.
Solche pathologischen Fälle führen dazu, dass es nicht sinnvoll ist
von einer eindeutigen Faktorisierung zu sprechen.

\begin{definition}[Nullteiler, Integritätsbereich]
Nullteiler \emph{Nullteiler} in einem Ring $R$ sind Zahlen
$a,b\in R^*=R\setminus\{0\}$ mit $ab=0$.
Ein Ring $R$ ohne Nullteiler heisst \emph{nullteilerfrei} oder ein
\emph{Integritätsbereich}.
\index{Nullteiler}%
\index{nullteilerfrei}%
\index{Integritatsbereich@Integritätsbereich}%
\end{definition}

Der euklidische Algorithmus für ganze Zahlen verwendet, dass der Rest
bei der ganzzahligen Division kleiner als der Divisor ist.
Polynome kann man nicht vergleichen.
Der euklidische Algorithmus für Polynome braucht aber, dass der Grad
des Restes kleiner wird als der Grad des Divisors.
Der Algorithmus ist also auf eine allgemeinere Situation übertragbar,
wenn ein Eigenschaft ähnlich dem Grad gefunden werden kann, die bei
der Division kleiner wird.

\begin{definition}[Euklidischer Ring]
\label{buch:ratint:euklid:def:euklidischerring}
Ein nullteilerfreier Ring  $R$ heisst ein \emph{euklidischer Ring},
wenn es eine Funktion $d\colon R^*\to\mathbb{Z}$
\begin{enumerate}
\item $d(r) \ge 0$ für alle $r\in R^*$.
\item $d(r)\le d(rs)$ für alle $r,s\in R^*$.
\item Sind $a,b\in R^*$, dann gibt es $q,r$ mit $a=qb+r$ mit $r=0$ oder
$d(r)<d(b)$.
\end{enumerate}
Die Funktion $d$ heisst auch die \emph{Gradfunktion} des euklidischen
Rings.
\index{Gradfunktion}%
Die letzte Eigenschaft kann etwas vereinfacht werden, wenn man
$d(0)=-\infty$ definiert.
\end{definition}

Für die ganzen Zahlen ist die $d\colon\mathbb{Z}^*\to\mathbb{Z}: a\to a$
die Funktion, die die ganzen Zahlen zu einem euklidischen Ring macht.
Für Polynome mit Koeffizienten in einem Körper ist es der Grad
$\deg\colon K[X] \to \mathbb{Z} : p(X) \mapsto \deg p$.

Ein weiteres interessantes Beispiel für einen euklidischen Ring ist
der Ring der ganzen gaussschen Zahlen.

\begin{beispiel}
Wir betrachten die Menge
\[
\mathbb{Z}[i]
=
\{a+bi\mid a,b\in\mathbb{Z}\}\subset\mathbb{C}.
\]
Sie heisst der \emph{Ring der ganzen gaussschen Zahlen}.
\index{gausssche Zahlen}%
Als Teilmenge des Körpers der komplexen Zahlen ist $\mathbb{Z}[i]$
sicher ein nullteilerfreier Ring.
Als Gradfunktion kann das Betragsquadrat 
\[
d
\colon
\mathbb{Z}[i] \to \mathbb{Z}
:
(a+bi) \mapsto a^2+b^2
\]
verwendet werden.
Es müssen die drei Eigenschaften der 
Definition~\ref{buch:ratint:euklid:def:euklidischerring}
verifizieren.
\begin{enumerate}
\item
Für $a+bi\ne 0$ ist $d(a+bi)=a^2+b^2>0$.
\item
Für das Produkt gilt für $(a+bi)\ne 0$ ist $d(a+bi)\ge 1$
\[
d\bigl(
(a+bi)(c+di)
\bigr)
=
(a^2+b^2)(c^2+d^2)
=
d(a+bi)d(c+di)
\ge
d(c+di)
\]
für all $c+di\in \mathbb{Z}[i]$.
\item
TODO
\end{enumerate}
Dies zeigt, dass die ganzen gaussschen Zahlen einen euklidischen
Ring bilden.
Es ist daher sinnvoll, vom grössten gemeinsamen Teiler ganzer
gaussscher Zahlen zu sprechen und es ist möglich, diese mit dem
euklidischen Algorithmus zu finden.
\end{beispiel}

Die ganzschen gaussschen Zahlen können zum Beispiel dazu verwendet
werden, dann folgenden fermatschen Zweiquadratesatz zu beweisen
\cite[Satz 10.4]{buch:luneburg}.

\begin{satz}[Fermat]
Ist $p$ eine Primzahl mit $p\equiv 1\mod 4$, dann gibt es ganze
Zahlen $a,b\in\mathbb{Z}$ mit $p=a^2+b^2$.
\end{satz}

%
% Polynome mit Polynomkoeffizienten
%
\subsection{Polynome mit Polynomkoeffizienten
\label{buch:ratint:euklid:subsection:polypoly}}




%
% 3-hermite.tex
%
% (c) 2026 Prof Dr Andreas Müller
%
\section{Quadratfreie Faktorisierung und die Methode von Hermite
\label{buch:ratint:section:hermite}}
\kopfrechts{Quadratfreie Faktorisierung und die Methode von Hermite}%
Gemäss dem Plan am Ende von Abschnitt~\ref{buch:ratint:section:bernoulli}
ist, ist eine Faktorisierung des Nenners mit sehr speziellen
Teilbarkeitseigenschaften nötig.
Im vorangegangenen Abschnitt wurde mit dem euklidischen Algorithmus
das dazu geeignete Werzeug gefunden.
Dieses soll jetzt dazu eingesetzt werden, die Faktorisierung zu
finden und den Reduktionsschritt für die Integrale durchzuführen.

%
% Quadratfreie Faktorisierung
%
\subsection{Quadratfreie Faktorisierung}
Die quadratfreie Faktorisierung löst das Problem, den Nenner ohne
Bestimmung der Nullstellen so zu faktorisieren, dass eine
Partialbruchzerlegung möglich wird, mit der die Stammfunktion
bestimmt werden kann.

%
% Mehrfache Nullstellen
%
\subsubsection{Mehrfache Nullstellen}
Sei $q(z) \in \mathbb{C}[z]$ ein Polynom mit einer $n$-fachen Nullstelle
$\alpha$.
Dies bedeutet, dass $q(z)$ in der Form
\[
q(z)
=
(z-\alpha)^n\cdot s(z)
\]
mit einem Polynom $s(z)\in\mathbb{C}[z]$ geschrieben werden kann.
Das Polynom $s(z)$ hat $\alpha$ nicht als Nullstelle, $s(\alpha)\ne 0$,
es ist also nicht durch $z-\alpha$ teilbar.

Nach der Produktregel ist die Ableitung von $q(z)$
\[
q'(z)
=
n
(z-\alpha)^{n-1}
\cdot s(z)
+
(z-\alpha)^n\cdot s'(z)
=
(z-\alpha)^{n-1} 
\cdot (s(z)+(z-\alpha)\cdot s'(z)).
\]
Für $z=\alpha$ ist der Klammerfaktor
\[
s(\alpha) + (\alpha-\alpha)\cdot s'(\alpha)
=
s(\alpha)
\ne 0,
\]
der Klammerfaktor hat $\alpha$ nicht als Nullstelle.
Für eine mehrfache Nullstelle ist $n>1$ und es folgt, dass $q'(z)$
die Zahl $\alpha$ als $n-1$-fache Nullstelle hat.

Diese Beobachtung lässt sich auch verallgemeinern.
Hat $q(z)$ einen Faktor $p(z)$, der $n$-fach in $q(z)$ vorkommt,
dann kann man
\[
q(z)
=
p(z)^n\cdot s(z)
\]
schreiben, wobei $s(z)$ nicht durch $p(z)$ teilbar ist.
Die Ableitung von $q'(z)$ ist dann nach der Produktregel
\[
q'(z)
=
n p(z)^{n-1} \cdot s(z) + p(z)^n \cdot s'(z)
=
q(z)^{n-1}\cdot (ns(z) + p(z) s'(z)).
\]
Der Faktor auf der rechten Seite ist nicht durch $p(z)$ teilbar,
weil $s(z)$ nicht durch $p(z)$ teilbar ist.
Es folgt, dass $q'(z)$ durch $p(z)^{n-1}$ teilbar ist.

Wenn also $q(z)$ einen mehrfachen vorkommenden Faktor hat,
dann ist dieser Faktor ein gemeinsamer Teiler von $q(z)$
und $q'(z)$.

\begin{beispiel}
TODO
\end{beispiel}

%
% Quadratfreie Polynome
%
\subsubsection{Quadratfreie Polynome}
Damit $q(z)$ und $q'(z)$ keinen nichtrivialen gemeinsamen Teiler
haben, darf jeder Linearfaktor nur ein einziges Mal vorkommen.
Solche Polynome heissen quadratfrei im Sinne der folgenden
Definition.

\begin{definition}
Ein Polynom in $\mathbb{C}[z]$ heisst \emph{quadratfrei}, wenn es
keine mehrfachen Nullstellen hat oder wenn es keine mehrfachen
Faktoren hat.
\index{quadratfrei}%
\end{definition}

Über den komplexen Zahlen zerfällt also ein quadratfreies Polynom
immer in ein Produkt von \emph{verschiedenen} Linearfaktoren.

Wenn ein Polynome $q(z)$ mehrfache Nullstellen hat, dann hat die
Ableitung die gleichen Nullstellen, aber mit einer um $1$ kleineren
Ordnung.
Der grösste gemeinsame Teiler von $q(x()$ und $q'(x)$ besteht
daher alle Faktoren, die mehr als einmal vorkommen, mit einer
gegenüber $q(z)$ um $1$ verminderten Vielfachheit.
Im Quotienten
\[
p(z)
=
\frac{q(z)}{\operatorname{ggT}(q(z),q'(z))}
\]
bleiben dann nur noch alle Linearfaktoren genau einmal stehen.

\begin{beispiel}
Das Polynom 
\[
q(z)
=
z^5-z^4-2z^3+2z^2+z-1
\]
hat die Ableitung
\[
q'(z)
=
5z^4-4z^3-6z^2+4z+1.
\]
Der euklidische Algoritmus liefert als grössten gemeinsamen
Teiler das Polnom
\[
g(z)
=
\operatorname{ggT}(q(z), q'(z))
=
z^3-z^2-1+1.
\]
Der Quotient
\[
\frac{q(z)}{g(z)}
=
z^2-1
\]
ist ein quadratfreies Polynom, welches alle Linearfaktoren von $q(z)$
enthält.
Es ist in diesem Fall einfach zu sehen, dass $z^2-1=(z-1)(z+1)$,
d.~h.~die einzigen Nullstellen von $q(x)$ sind $\pm 1$.

Man kann aber die Faktoren auch noch auf etwas systematischere Art
finden.
Dazu berechnet man den grössten gemeinsamen Teiler von $g(z)$ und $g'(z)$,
wegen $g'(z) = 3z^2 -2z -1$ ist
ist
\[
\operatorname{ggT}(g(z),g'(z))
=
z-1
\]
und
\[
\frac{g(z)}{\operatorname{ggT}(g(z),g'(z))}
=
z^2-1.
\]
Daraus kann man schliessen, dass es zwei Linearfaktoren von $q(z)$
gibt.
Einer davon kommt genau zweimal vor, der andere ist $z-1$ und kommt
dreimal vor.
Durch Division $(z^2-1)/(z-1)=z+1$ findet man jetzt, dass $z+1$
der Faktor ist, der genau zweimal ist.
Wir haben somit die Faktorisierung
\[
q(z)
=
(z+1)^2 (z-1)^3
\]
gefunden.
\end{beispiel}

Das Beispiel zeigt, dass die Suche nach quadratfreien Faktoren
mithilfe des grössten gemeinsamen Teilers manchmal sogar eine
Faktorisierung des Polynoms liefern kann.

%
% Die Aufgabe der Faktorisierung in quadratfreie Faktoren
%
\subsubsection{Die Aufgabe der Faktorisierung in quadratfreie Faktoren}
Ein beliebiges normiertes Polynom $q(z)$ kann einzelne Nullstellen
mehrfach enthalten.
Man kann aber immer noch eine Faktorisierung der Form
\[
q(z)
=
\prod_{k=1}^n
(z-\alpha_k)^{n_k},
\]
wobei alle Nullstellen $\alpha_k$ verschieden sind und $n_k$ ihre
Vielfachheit ist.
Fasst man die Faktoren mit gleicher Vielfachheit zusammen, kann man
zu jeder Vielfachheit $l$ das Polynom
\[
q_l(z) = \prod_{n_k=l} (z-\alpha_k)
\]
bilden.
Da die Faktoren von $q_l(z)$ im Polynom $q(z)$ genau $l$-mal vorkommen,
kann man das Polynom $q(z)$ aus
\[
q(z)
=
q_1(z)
\cdot
q_2(z)^2
\cdot
q_3(z)^3
\cdot
\ldots
\cdot
q_m(z)^m
\]
rekonstruieren.
Nach Konstruktion kann ein Linearfaktor nur immer in genau einem
der Faktoren $q_l(z)$ vorkommen, die Faktoren sind also paarweise
teilerfremd.

\begin{definition}[quadratfreie Faktorisierung]
\index{quadratfreie Faktorisierung}%
\index{Faktorisierung, quadratfrei}%
Die \emph{quadratfreie Faktorisierung} eines Polynoms $q(z)$ ist eine
Faktorisierung der Form
\[
q(z)
=
a
\prod_{l=1}^m q_l(z)^l
\]
mit teilerfremden Faktoren.
\end{definition}

Die Bestimmung der quadratfreien Faktorisierung aus der Zerlegung
in Linearfaktoren zeigt zwar, dass es eine solche Faktorisierung
gibt, ist aber kein praktikabler Weg, sie algorithmisch zu bestimmen.
daher soll in diesem Abschnitt soll die folgende Aufgabe gelöst werden.

\begin{aufgabe}
Formuliere einen Algorithmus, der zu einem beliebigen rationalen Polynom
$q(z)\in\mathbb{Q}[z]$ die quadratfreie Faktorisierung findet.
\end{aufgabe}

%
% Berechnung der quadratfreien Faktorisierung
%
\subsubsection{Berechnung der quadratfreien Faktorisierung}
Aus dem Beispiel lesen wir ab, dass der grösste gemeinsame
Teiler $g(z)=\operatorname{ggT}\bigl(q(z),q'(z)\bigr)$ von $q(z)$
und $q'(z)$ alle mehrfachen Faktoren enthält und der Quotient
$q(z)/g(z)$ alle Linearfaktoren genau einmal.
Diese Eigenschaft soll jetzt dazu verwendet werden, die Faktoren
$q_l(z)$ zu finden.
Dazu konstruieren wir zunächst zwei Folgen $g_k(z)$ und $p_k(z)$
wie folgt.
Das Polynom $g_0(z)$ ist $g_0(z)=q(z)$.
Die Polynome $g_k(z)$ und $p_k(z)$ mit $k>0$ sind rekursiv definiert
durch
\begin{align*}
g_{k+1}(z) &= \operatorname{ggT}\bigl(g_k(z),g'_k(z)\bigr)
&&\text{und}&
p_{k+1}(z) &= \frac{g_k(z)}{g_{k+1}(z)}.
\end{align*}
Nach Konstruktion sind die Polynom $p_k(z)$ alle quadratfrei.
Sie enthalten alle Linearfaktoren von $g_k(z)$.
Nach Konstruktion ist $g_{k+1}(z)$ immer ein echter Teiler von $g_k(z)$.
Der Grad nimmt daher in jedem Schritt ab.
Die Konstruktion bricht daher nach endlich vielen Schritten ab.

In jedem Iterationsschritt nimmt der Exponent der in $g_k(z)$
vertretenen Linearfaktoren um $1$ ab.
Das Ende $g_k(z)=1$ wird also genau in dem Moment erreicht, 
in dem der grösste Exponent auf $1$ reduziert worden ist.
Somit ist $k$ der grösste Exponent eines Linearfaktors.

Die Polynome $p_k(z)$ enthalten alle Linearfaktoren, die
mindestens $k$ mal in $q(z)$ aufreten.
Es folgt, dass $p_{k+1}(z)$ ein Teiler von $p_k(z)$ ist und dass
der Quotient $p_k(z)/p_{k+1}(z)$ aus den Faktoren besteht, die 
genau $k$-mal in $q(z)$ vorkommen.
Dies ist genau, was die Polynome $q_l(z)$ definiert.
Dies beweist den folgenden Satz.

\begin{satz}[quadratfreie Faktorisierung]
Sei $q(z)\in K[z]$ ein normiertes Polynom.
Setze $q_0(z)=q(z)$ und bilde rekursiv die Folgen
\begin{align*}
g_k(z) &= \operatorname{ggT}\bigl(g_{k-1}(z),g_{k-1}'(z)\bigr)
&
p_k(z) &= \frac{g_{k-1}(z)}{g_k(z)}
&
q_k(z) &= \frac{p_k(z)}{p_{k+1}(z)}.
\end{align*}
Sei $m$ dasjenige $k$, für das $p_k(z)=1$ ist.
Dann ist
\[
q(z)
=
\prod_{l=1}^m q_l(z)^l
\]
die quadratfreie Faktorisierung.
\end{satz}

\begin{beispiel}
Wir betrachten das Polynom
\begin{align*}
q(z)
&=
z^9+z^8-5z^7-5z^6+9z^5+9z^4-7z^3-7z^2+2z+2
\intertext{mit der Ableitung}
q'(z)
&=
9z^8+8z^7-35z^6-30z^5+45z^4+36z^3-21z^2-14z+2
\end{align*}
Die Polynome $g_k(z)$, $p_k(z)$ und $q_l(z)$ sind
\begin{align*}
g_1(z)
&=
z^5+z^4-2z^3-2z^2+z+1,
&
p_1(z)
&=
z^4-3z^2+2,
&
q_1(z)
&=
z^2-2,
\\
g_2(z)
&=
z^3+z^2-z-1,
&
p_2(z)
&=
z^2-1,
&
q_2(z)
&=
1,
\\
g_3(z)
&=,
z+1
&
p_3(z)
&=
z^2-1,
&
q_3(z)
&=
z-1,
\\
g_4(z)
&=
1,
&
p_4(z)
&=
z+1,
&
q_4(z)
&=
z+1,
\\
g_5(z)
&=
1,
&
p_5(z)
&=
1.
&
&
\end{align*}
Tatsächlich kann man nachrechnen, dass
\[
q(z)
=
(z^2-2)(z-1)^3(z+1)^4
\]
ist.
Man beachte aber, dass diese Faktorisierung mit $z^2-2$ auch einen Faktor
enthält, der kein Linearfaktor ist.
Da $z^2-2 = (z+\!\sqrt{2})(z-\!\sqrt{2})$, gibt es über $\mathbb{Q}$
gar keine Faktorisierung in Linearfaktoren.
\end{beispiel}


%
% Reduktion des Integrals
%
\subsection{Reduktion des Integrals}
Mit dem Werkzeug der quadratfreien Faktorisierung lässen sich jetzt
die für den rationalen Teil des Integrals nötigen Schritte des Plans
vom Ende von Abschnitt~\ref{buch:ratint:section:bernoulli} durchführen.
Zu diesem Zwck muss erst eine Partialbruchzerlegung gefunden werden.
Anschliessend müssen die einzelnen Terme der Partialbruchzerlegung
integriert werden.

%
% Partialbruchzerlegung zur quadratfreien Faktorisierung
%
\subsubsection{Partialbruchzerlegung zur quadratfreien Faktorisierung}
Der wesentliche Schritt zur Konstruktion der Partialbruchzerlegung
ist der folgende Satz.

\begin{satz}
Sei 
\[
f(z)
=
\frac{p(z)}{q(z)^m r(z)},
\]
wobei $q(z)$ teilerfremd zu $r(z)$ ist.
Dann gibt es 
\[
f(z)
=
f_0(z)
+
\frac{s(z)}{r(z)}
+
\sum_{k=1}^m \frac{t_k(z)}{q(z)^k},
\]
wobei $f_0(z)$ ein Polynom ist, $s(z)$ ein Polynom mit Grad
$\deg s(z)<\deg r(z)$ und die Polynome $t_k(z)$ Grad $\deg t_k(z)<\deg q(z)$
haben.
\end{satz}

\begin{proof}
Da $q(z)$ und $r(z)$ teilerfremd sind, sind auch $q(z)^m$ und $r(z)$
teilerfremd.
Daher gibt es $s(z)$ und $t(z)$ mit
\begin{equation}
p(z)
=
s(z)q(z)^m + t(z)r(z).
\label{buch:ratint:hermite:eqn:p=sq+tr}
\end{equation}
Einsetzen in den Bruch ergibt
\begin{align}
f(z)
&=
\frac{p(z)}{q(z)^m r(z)}
=
\frac{s(z)q(z)^m + t(z)r(z)}{q(z)^mr(z)}
=
\frac{s(z)}{r(z)}
+
\frac{t(z)}{q(z)^m}.
\label{buch:ratint:hermite:eqn:srtq}
\end{align}
Falls der Zähler einen Grad mindestens so gross wie der Nenner
hat, kann durch Division des Zählers mit Rest das Polynom $f_0(z)$
abgespaltet werden.

Es muss jetzt nur noch gezeigt werden, dass das der zweite Bruch
\eqref{buch:ratint:hermite:eqn:srtq}
in eine Summe der verlangten Art umgefomrt werden kann.
Dazu schreibt man den Zähler als
\begin{equation*}
t(z)
=
a_m(z) q(z) + t_{m}(z)
\end{equation*}
mit $\deg b(z) < \deg q(z)$.
Dann wird
\begin{align*}
\frac{t(z)}{q(z)^m}
&=
\frac{a(z)q(z)+b(z)}{q(z)^m}
=
\frac{a(z)}{q(z)^{m-1}}
+
\frac{b(z)}{q(z)^m}.
\intertext{Durch wiederholte Anwendung dieser Beobachtung auf
den neuen Zähler $a_m(z) = a_{m-1}(z)q(z)+t_{m-1}(z)$
wird daraus eine Summe der Form}
&=
\frac{t_1(z)}{q(z)}
+
\frac{t_2(z)}{q(z)^2}
+
\ldots
+
\frac{t_m(z)}{q(z)^m}.
\end{align*}
Damit ist der Satz vollständig bewiesen.
\end{proof}

Durch Anwendung auf eine Faktorisierung von $r(z)$ kann schrittweise
die Partialbruchzerlegung mit den Teilnennern $q_l(z)^k$, $k\le l$
für alle $l$ gefunden werden.
Man beachte, dass dazu wieder nur die Polynomdivision
und für die Identität 
\eqref{buch:ratint:hermite:eqn:p=sq+tr}
der euklidische Algorithmus benötigt wird.
Beide können für Polynome in einer beliebigen Körpererweiterung
für den Koeffizientenkörper durchgeführt werden.


%
% Integration des rationalen Teils
%
\subsubsection{Integration des rationalen Teils}
Ganz analog der Vorgehensweise in der Methode von Bernoulli im
Abschnitt~\ref{buch:ratint:section:bernoulli} können auch die
Integranden mit Potenzen quadratfreier Polynome $q(z)$ im Nenner
auf den Nenner $q(z)$ reduziert werden.

\begin{satz}
\label{buch:ratint:hermite:satz:reduktion}
Sei $q(z)$ ein quadratfreies Polynom und $p(z)$ ein beliebiges
Polynom von kleinerem Grad $\deg q(z)<\deg p(z)$.
Dann gibt es Polynome $s(z)$ und $t(z)$ Polynome mit
\begin{equation}
s(z)\cdot q(z)
+
t(z)\cdot q'(z)
=
p(z),
\label{buch:ratint:hermite:ggtsz}
\end{equation}
für die zudem
\begin{align*}
\int
\frac{p(z)}{q(z)^m}
\,dz
&=
\int
\frac{s(z)+t'(z)/(m-1)}{q(z)^{m-1}}
\,dz
+
\frac{t(z)}{(1-m)q(z)^{m-1}}
\end{align*}
wird.
\end{satz}

\begin{proof}
Da $q(z)$ quadratfrei ist sind  $q(z)$ und $q'(z)$ teilerfremd.
Nach dem euklidischen Algorithmus für Polynome gibt es daher
Polynome $s(z)$ und $t(z)$ derart, dass \eqref{buch:ratint:hermite:ggtsz}
gilt.
Einsetzen ins Integral ergibt
\begin{align*}
\int
\frac{p(z)}{q(z)^m}
\,dz
&=
\int
\frac{s(z)\cdot q(z)+t(z)\cdot q'(z)}{q(z)^m}
\,dz
\\
&=
\int
\frac{s(z)}{q(z)^{m-1}}
\,dz
+
\int
t(z)
\cdot
\frac{q'(z)}{q(z)^m}
\,dz
\intertext{Auf das zweite Integral kann partielle Integration
angewendet werden.
Das Polynom $p(z)$  wird abgeleitet während der Bruch
$q'(z)/q(z)^m$ integriert wird.
So entsteht}
&=
\int
\frac{s(z)}{q(z)^{m-1}}
\,dz
+
t(z)
\frac{1}{(1-m)q(z)^{m-1}}
-
\int
\frac{t'(z)}{(1-m)q(z)^{m-1}}
\,dz
\\
&=
\int
\frac{s(z)+t'(z)/(m-1)}{q(z)^{m-1}}
\,dz
+
t(z)
\frac{1}{(1-m)q(z)^{m-1}}.
\end{align*}
Damit ist der Satz bewiesen.
\end{proof}

Durch wiederholte Anwendung des
Satzes~\ref{buch:ratint:hermite:satz:reduktion}
kann man die Potenz des Nenners immer weiter reduzieren, bis
nur noch das Integral
\[
\int
\frac{p(z)}{q(z)}
\,dz
\]
bestimmt werden muss.
Dieser Teil des Integrationsproblems heisst auch der 
\emph{logarithmische Teil} des Integrals.
\index{logarithmischer Teil}%

%
% Die Methode von Horowitz-Ostrogradski
%
\subsection{Die Methode von Horowitz-Ostrogradski}
Die Berechnung der quadratfreien Zerlegung sowie die Reduktion
der Nennerpotenz mit der Methode von Hermite ist zwar algebraisch
einfach, verlangt aber die häufig wiederholte Berechnung von 
grössten gemeinsamen Teilern von Polynomen mit dem erweiterten
euklidischen Algorithmus.
Da jetzt aber die erwartete Form des rationalen und des logarithmischen
kennen, können die zugehörigen Zählerpolynome auch direkt mithilfe
eines linearen Gleichungssystems gesucht werden.
Dies ist die Methode von Horowitz, die allerdings schon früher von
Ostrogradski gefunden wurde.



%
% 4-rothsteintrager.tex
%
% (c) 2026 Prof Dr Andreas Müller
%
\section{Die Methode von Rothstein-Trager
\label{buch:ratint:section:rothsteintrager}}
\kopfrechts{Die Methode von Rothstein-Trager}%
Die Methode von Rothstein-Trager ermöglicht, auch den logarithmischen
Teil
\begin{equation}
I
=
\int
\frac{p(z)}{q(z)}
\,dz
\label{buch:ratint:rothsteintrager:eqn:integral}
\end{equation}
für einen quadratfreien Nenner $q(z)$ und einen dazu teilerfremden
Zähler $p(z)$ mit $\deg p(z)<\deg q(z)$ zu berechnen.

%
% Residuen und die Form des Integrals
%
\subsection{Residuen und die Form des Integrals
\label{buch:ratint:rothsteintrager:subsection:residuen}}
In den komplexen zahlen lässt sich der quadratfreie Nenner $q(z)$ von
\eqref{buch:ratint:rothsteintrager:eqn:integral}
in ein Produkt verschiedner Linearfaktoren
\begin{equation}
q(z)
=
(z-\alpha_1)\cdots(z-\alpha_n)
\label{buch:ratint:rothsteintrager:eqn:faktorisierung}
\end{equation}
zerlegen.
Da die Nullstellen alle verschieden sind, hat die Partialbruchzerlegung
des Integranden die Form
\begin{equation}
\frac{p(z)}{q(z)}
=
\frac{A_1}{z-\alpha_1}
+
\ldots
+
\frac{A_n}{z-\alpha_n}.
\label{buch:ratint:rothsteintrager:eqn:integrand}
\end{equation}
In dieser Form lässt sich das Integral jedoch sofort angeben, es ist
\begin{equation}
\int
\frac{p(z)}{q(z)}
\,dz
=
A_1\log(z-\alpha_1)
+
\ldots
+
A_n\log(z-\alpha_n).
\label{buch:ratint:rothsteintrager:eqn:integralform}
\end{equation}
Der Ausdruck kann noch etwas vereinfacht werden, wenn zwei der
Koeffizienten gleich sind.
Ist nämlich $A_k=A_l$, dann können die beiden Summanden in
$A_k\log\bigl((z-\alpha_k)(z-\alpha_k)\bigr)$
zusammengefasst werden.
Für entgegengesetzte Koeffizienten $A_k=-A_l$ folgt entsprechend
\[
A_k\log \frac{z-\alpha_k}{z-\alpha_l}.
\]
Diese Form tritt zum Beispiel bei der Integration reeller Funktionen
auf, wo die Nullstellen des Nenners paarweise konjugiert komplex sind.

Wenn die Faktorisierung
\eqref{buch:ratint:rothsteintrager:eqn:faktorisierung}
von $q(z)$ bekannt ist, ist das Integrationsproblem an dieser
Stelle gelöst.
Wie früher bereits dargelegt wurde, kann man jedoch nicht davon
ausgehen, dass das Faktorisierungsproblem gelöst werden kann.
Es müssen also alternative Methoden gesucht werden, die ohne die 
die Faktorisierung auskommen.

%
% Die Koeffizienten A_k und Residuen
%
\subsubsection{Die Koeffizienten $A_k$ und Residuen}
Die Form 
\eqref{buch:ratint:rothsteintrager:eqn:integrand}
des Integranden des logarithmischen Teils zeigt, dass er isolierte
Singularitäten hat und daher um jede Singularität in eine Laurent-Reihe
entwickelt werden kann.
Man kann sogar sagen, dass die Singularitäten Pole erster Ordnung
sein müssen.

Sie jetzt $\alpha$ eine dieser Pole, dann kann der Integrand
in eine Laurent-Reihe
\[
\frac{p(z)}{q(z)}
=
\frac{b_1}{z-\alpha}
+
a_0+a_1(z-\alpha) + a_2(z-\alpha)^2+\ldots
\]
um den Punkt $\alpha$ entwickelt werden.
Gliedweise Integration der Reiher ergibt in einer Umgebung von $\alpha$
\[
\int
\frac{p(z)}{q(z)}
\,dz
=
b_1\log(z-\alpha)
+
a_0(z-\alpha) + \frac{a_1}{2}(z-\alpha)^2 + \frac{a_2}{3}(z-\alpha)^3
+\ldots
\]
Durch Vergleich mit
\eqref{buch:ratint:rothsteintrager:eqn:integralform}
folgt, dass die gesuchten Koeffizienten $A_k$ die Residuen
des Integranden an den Nullstellen von $q(z)$ sein müssen.
Diese Koeffizienten lassen sich aber direkt aus den Polynomen 
$p(z)$ und $q(z)$ berechnen.

\begin{satz}
\label{buch:ratint:rothsteintrager:satz:residuum}
Sind $p(z)$ und $q(z)$ teilerfremd und $q(z)$ quadratfrei und
$\alpha$ eine Nullstelle des Nenners.
Dann ist das Residuum von $p(z)/q(z)$ an der Stelle $\alpha$
\[
A
=
\operatorname{res}\biggl(\alpha;\frac{p(z)}{q(z)}\biggr)
=
\frac{p(\alpha)}{q'(\alpha)}.
\]
\end{satz}

Der Satz ist eine Folge des folgenden, algemeineren Satzes.
Er gilt für einen beliebigen Bruch von holomorphen Funktionen,
dessen Nenner eine einfache Nullstelle hat.

\begin{satz}
Ist $f(z)$ in einer Umgebung von $\alpha$ holomorph und hat $g(z)$
eine Nullstelle erster Ordnung in $\alpha$, dann ist das Residuum
von $f(z)/g(z)$ an der Stelle $\alpha$ der Wert $f(\alpha)/g'(\alpha)$.
\end{satz}

\begin{proof}
Da $g(z)$ eine Nullstelle erster Ordnung an der Stelle $\alpha$
hat, kann man
\[
g(z) = (z-\alpha)\bigl(g'(\alpha)+ r(z)(z-\alpha)\bigr)
\]
und damit
\[
\frac{f(z)}{g(z)}
=
\frac{f(z)}{(z-\alpha)\bigl(g'(\alpha)+r(z)(z-\alpha)\bigr)}
=
\frac{1}{z-\alpha}
\cdot
\frac{f(z)}{g'(\alpha)+r(z)(z-\alpha)}
=
\frac{1}{z-\alpha}
\cdot
h(z)
\]
mit einer in einer Umgebung von $\alpha$ holomorphen Funktion
$h(z)$ schreiben.
Nach Satz~\ref{buch:analytisch:laurent:satz:laurentreihe} ist
das Residuum das Integral
\begin{align*}
\frac{1}{2\pi i}
\oint_\gamma
\frac{f(z)}{g(z)}
\,dz
&=
\frac{1}{2\pi i}
\oint_\gamma
\frac{1}{z-\alpha}
\frac{f(z)}{g'(\alpha)+r(z)(z-\alpha)}
\,dz
\\
&=
\frac{1}{2\pi i}
\oint_\gamma
\frac1{z-\alpha}
\cdot
h(z)
\,dz.
\intertext{Nach dem Integralsatz~\ref{buch:integration:satz:residuensatz}
ergibt das Integral}
&=
h(\alpha)
=
\frac{f(\alpha)}{g'(\alpha)}.
\end{align*}
Dies beweist den Satz.
\end{proof}

%
% Der Satz von Rothstein-Trager
%
\subsection{Der Satz von Rothstein-Trager}
Die Form des logarithmischen Teils des Integrals ist zwar bekannt,
und der Satz~\ref{buch:ratint:rothsteintrager:satz:residuum}
ermöglicht eine direkte Berechnung des Koeffizienten, sobald
die Nullstellen bekannt sind, aber genau die Bestimmung der
Nullstellen war ja der Punkt, warum die Methode von Bernoulli-Methode
als unzureichend erkannt wurde.
Dies ändert der folgende Satz von Rothstein-Trager.

\begin{satz}[Rothstein-Trager]
\label{buch:ratint:rothsteintrager:satz:rothsteintrager}
Ist $p(z)$ teilerfremd zum quadratfreien Polynom $q(z)$ mit
$\deg p(z)<\deg q(z)$, dann ist
\[
\int
\frac{p(z)}{q(z)}
\,dz
=
\sum_{i=1}^m c_i \log v_i(z),
\]
wobei die $c_i$ die möglicherweise komplexen Zahlen $c$ sind, derart,
dass $p(z)-cq'(z)$ und $q(z)$ einen gemeinsamen Faktor haben.
Die zugehörigen $r_i(z)$ sind
\[
r_i(z)
=
\operatorname{ggT}\bigl(p(z)-c_iq'(z),q(z)\bigr).
\]
\end{satz}

\begin{proof}
TODO
\end{proof}

Wenn es gelingt, eine algebraische Methode zur Bestimmung der
Koeffizienten $c_i$ zu finden, können auch die Faktoren 
$r_i(z)$ gefunden werden.
Man beachte, dass nicht mehr eine Faktorisierung von $q(z)$
verlangt wird, sondern nur noch die Möglichkeit, gemeinsame
Faktoren zu finden, was im Prinzip mit dem euklidischen Algorithmus
möglich ist.

%
% Gemeinsame Nullstellen und die Resultante
%
\subsubsection{Gemeinsame Nullstellen und die Resultante}
Der Satz~\ref{buch:ratint:rothsteintrager:satz:rothsteintrager}
verlagert das Problem, die Nullstellen des Nenners zu finden darauf,
die Koeffizienten $c_i$ zu finden.
Dies ist aus zwei Gründen einfacher:
\begin{enumerate}
\item
Oft sind einige der Koeffizienten $A_k$ gleich.
In diesen Fällen ist es nicht nötig, die zugehörigen Linearfaktoren
zu kennen, es genügt das Produkt dieser Linearfaktoren zu finden,
dies ist das Polynom $r_i(z)$ im
Satz~\ref{buch:ratint:rothsteintrager:satz:rothsteintrager}.
\item
Für das Kriterium, einen gemeinsamen Teiler zu finden gibt es
ein einfaches algebraisches Kriterium, nämlich die Resultante.
\end{enumerate}

In Kapitel~\ref{chapter:resultante} wird beschrieben, wie die 
Resultante des Problem, die Koeffizienten $c_i$ zu finden, lösen
kann.




%\uebungsabschnitt
%
%\aufgabetoplevel{chapters/060-ratint/uebungsaufgaben}
%\begin{uebungsaufgaben}
%\uebungsaufgabe{601}
%\uebungsaufgabe{602}
%\end{uebungsaufgaben}
%\enduebungsabschnitt

